By Thm.~\ref{thm:intconv-sound}, IntConv is a sound sufficient condition for ExtConv. A natural question is: for what circuits is IntConv \emph{complete}? That is, when does mathematical convergence according to ExtConv also imply IntConv? For such circuits, a system following the IntConv criterion faithfully implements the theoretical specification. Moreover, because IntConv and ExtConv coincide on these circuits, a sound and complete ExtConv detector also detects IntConv. This permits specialized detectors such as FirstZero for the naive and semi-naive circuits instead of the general StFP detector from \secref{sec:detector}.

To answer this question, we introduce a key property called \emph{convergence-complete} (Conv-complete) circuits. They are just the circuits mentioned above for which IntConv is complete. In \secref{sec:regular-circuits}, we use \emph{regular circuits} as a name for the fragment generated by the standard DBSP constructs for lifted scalar computation, composition, Datalog recursion, and while recursion. The name lets us state a uniform completeness result for this existing recursive sublanguage: all regular circuits are Conv-complete (Thm.~\ref{thm:regular-convcomplete}).
Moreover, in \secref{sec:inc-opt}, we prove that the DBSP's incremental optimization ($\mathsf{plift}$ and $\mathsf{incOpt}$) preserves Conv-completeness (Thm.~\ref{thm:pushLifting-convcomplete} and Thm.~\ref{thm:incOpt-convcomplete}). As a result, all regular circuits and their incrementally optimized versions are Conv-complete.

Along with Conv-completeness, its local version FP-completeness is important for us to establish these results. We define them as follows:
\begin{definition}[Conv-complete Circuit]
A circuit $c$ is \emph{Conv-complete} iff for every input $x$,
\[
\text{ExtConv}(c, x) \implies \text{IntConv}(c, x).
\]
\end{definition}
\begin{definition}[FP-complete Circuit]
Let $c$ be a well-formed circuit.
We say $c$ is \emph{FP-complete} if:
\begin{itemize}
    \item \textbf{Level-1 completeness:}
    for every input $x$ and index $n$, if $\text{ExtConv}(c, x)$ and $\text{ExtFP}_1(c, x, n)$ hold, then there exists $n' \in \N$ such that $\text{IntFP}_1(c, x, n')$ holds.
    \item \textbf{Level-2 completeness (for level-2 circuits):}
    for every nested stream input $x$ and stream $\mathbf{b}$, if $\text{ExtConv}(c, x)$ and $\text{ExtFP}_2(c, x, \mathbf{b})$ hold, then there exists a stream $\mathbf{b}'$ such that $\text{IntFP}_2(c, x, \mathbf{b}')$ holds.
\end{itemize}
\end{definition}

We have established some compositional properties of Conv-completeness and FP-completeness. In particular, Conv-completeness is closed under all circuit constructs except for $\mathsf{bracket}$, while FP-completeness is closed under all circuit constructs except for $\mathsf{bracket}$ and sequential composition $\mathsf{seq}$. These results are useful for us to establish the important results in this section.


% \subsection{Compositional Closure Properties of FP-completeness and Conv-completeness}\label{sec:complete-algebra}
% In short, FP-complete circuits and Conv-complete circuits are closed under most circuit constructs, but there are notable exceptions.

% \begin{lemma}[Compositional Closure Properties of Conv-complete Circuits]
%     All individual nodes (including primitive, standard and temporal nodes) are Conv-complete. Conv-complete circuits are closed under all other circuit constructs except for $\mathsf{bracket}$. 
% \end{lemma}

% \begin{lemma}[Bracket Rule for Conv-complete Circuits]
%     If $c$ is Conv-complete and FP-complete, then $\mathsf{bracket}(c)$ is Conv-complete.
% \end{lemma}

% These results immediately imply the following lemma. It means that we can compose FP-complete sub-circuits to get Conv-complete circuits, bridging the local fixpoint guarantee and the global convergence guarantee.

% \begin{lemma}[Conv-completeness from FP-completeness]\label{thm:convcomplete-global-criterion}
% If every sub-circuit of $c$ in the form $\mathsf{bracket}(c_0)$ has FP-complete inner circuit $c_0$, then $c$ is Conv-complete.
% \end{lemma}

% Compared to Conv-completeness, FP-completeness has more subtle closure properties.
% \begin{lemma}[Compositional Closure Properties of FP-complete Circuits]
%     All individual nodes are FP-complete. FP-complete circuits are closed under all other circuit constructs except for $\mathsf{bracket}$ and sequential composition $\mathsf{seq}$. 
% \end{lemma}  

% The most surprising result is that FP-complete circuits are not closed under sequential composition. To see this, consider the circuit $\mathsf{I} \seq \mathsf{D}$ on integers. Both $\mathsf{I}$ and $\mathsf{D}$ are FP-complete, but their sequential composition is not: when input is a constant stream of $1$, its output is also a constant stream of $1$, but the internal stream between $\mathsf{I}$ and $\mathsf{D}$ is $\lambda n.\, n+1$, which never reaches a fixpoint.

% This counterexample gives us an important insight on the hardness of the FPD problem. We can see that sequential composition hides the information of internal streams, allowing internally unstable circuits to behave stably externally. This may be one of the root causes of the FPD problem and the reason why IntFP is able to solve this problem.
% Nevertheless, for some good circuits, FP-completeness can be preserved by $\mathsf{seq}$ and $\mathsf{bracket}$, as the following lemma shows. It will be important for proofs in the next section.
% \begin{lemma}\label{thm:lifted-scalar-FPcomplete}
%     For a level-1 circuit $c_0$, if $c_0$ is FP-complete and $\means{c_0} = \lift{f}$ for some scalar function $f$, then $\mathsf{bracket}(c_0)$ is FP-complete. Additionally, if $c_1$ is FP-complete, then $c_0 \seq c_1$ is FP-complete.
% \end{lemma}

\subsection{Regular Circuits Are Conv-complete}\label{sec:regular-circuits}

Recall that we can use the DBSP circuit shown in Fig.~\ref{fig:naive-dbsp} to express a Datalog-style recursive query. We can express it in our formal language as follows.
\begin{definition}[Datalog Recursive Query]
    For a level-1 circuit $c_0$ that represents a lifted scalar query, the query body (the circuit between $\delta_0$ and $\elim$) of a Datalog-style recursive query is (the same as Eqn.~(\ref{eq:naive-body}))
    \[\mathsf{datalogBody}(c_0) := \Delta(\mathsf{loop}(c_0))\]
    and the streaming query (the lifted version of the whole circuit) is
    \[\mathsf{datalogQuery}(c_0) := \mathsf{bracket}(\mathsf{lifting}(\mathsf{datalogBody}(c_0)))\]
\end{definition}

Beyond recursive relational queries, the following \emph{while} program is used to show the Turing-completeness of DBSP~\cite{budiu2025dbsp}, where $Q$ is an arbitrary scalar function:
{\footnotesize
\begin{lstlisting}[language=Pascal]
x := i
while (x != Q(x))
  x := Q(x)
\end{lstlisting}
}
which can be implemented in DBSP as:
\begin{center}
\begin{tikzpicture}[>=latex]
    \node[] (input) {$i$};
    \node[block, right of=input] (delta) {$\delta_0$};
    \node[block, circle, right of=delta, inner sep=0cm] (p) {$+$};
    \node[block, right of=p] (Q) {$\lift Q$};
    \node[block, right of=Q] (D) {$\D$};
    \node[block, right of=D] (S) {$\elim$};
    \node[right of=S] (output)  {$x$};
    \node[block, below of=p, node distance=.7cm] (z) {$\zm$};
    \draw[->] (input) -- (delta);
    \draw[->] (delta) -- (p);
    \draw[->] (p) -- (Q);
    \draw[->] (Q) -- node (mid) {} (D);
    \draw[->] (D) -- (S);
    \draw[->] (mid.center) |- (z);
    \draw[->] (S) -- (output);
    \draw[->] (z) -- (p);
\end{tikzpicture}
\end{center}
In our formal language:
\begin{definition}[While Loop Programs]
    For a level-1 circuit $c_0$ representing a lifted scalar function, the query body of a while loop program is
    \[\mathsf{whileBody}(c_0) := \mathsf{loop}(\mathsf{add} \seq c_0) \seq \mathsf{D}\]
    and the streaming query is
    \[\mathsf{whileQuery}(c_0) := \mathsf{bracket}(\mathsf{lifting}(\mathsf{whileBody}(c_0)))\]
\end{definition}

Both query constructors above belong to the standard DBSP recursive sublanguage. Their results are level-1 circuits representing lifted scalar functions, so they can be composed or placed inside further recursive computations to form nested loops. We give the name \emph{regular circuits} to the sublanguage generated by these constructs so that we can state and prove FP- and Conv-completeness uniformly. We use \emph{lifted scalar circuit} for the base circuits that implement lifted scalar or relational queries.

\begin{definition}[Regular Circuit]
    \emph{Regular Circuits}, denoted by $\text{RegularCkt}$, is the class of level-1 circuits defined inductively as follows:
    \begin{itemize}
        \item \textbf{Lifted Scalar Circuit:} Any level-1 circuit composed (both sequentially and in parallel) of only primitive and standard nodes is a \emph{lifted scalar circuit}, which is also a regular circuit.
        \item \textbf{Sequence and Parallel Composition:} If $c_1$ and $c_2$ are regular circuits, their sequential composition $c_1 \seq c_2$ and parallel composition $c_1 \pll c_2$ are regular circuits.
        \item \textbf{Datalog Query:} Given a regular circuit $c_0$, $\mathsf{datalogQuery}(c_0)$ is a regular circuit.
        \item \textbf{While Loop:} Given a regular circuit $c_0$, $\mathsf{whileQuery}(c_0)$ is a regular circuit.
    \end{itemize}
\end{definition}

Using the inductive structure, it is easy to show that all regular circuits are FP-complete:
% Because all regular circuits represent a lifted scalar function, by leveraging Lemma~\ref{thm:lifted-scalar-FPcomplete},
%  it is easy to show that they are all FP-complete.
\begin{lemma}[Regular Circuits are FP-complete]
    Any regular circuit $c$ is FP-complete. In particular, for any $x, n$, $\text{ExtFP}_1(c, x, n) \implies \text{IntFP}_1(c, x, n)$.
\end{lemma}

We proceed to show that Datalog queries are Conv-complete. We first show that $\mathsf{datalogBody}$ satisfies a variant of FP-completeness tailored for a recursive query body, whose input must be in the form of $\delta_0(x)$ and the fixed output is zero.
\begin{lemma}[The iteration bound of $\mathsf{datalogBody}$]\label{thm:datalog-body-fpcomplete}
    If $c_0$ is a Conv-complete regular circuit, let $c := \mathsf{datalogBody}(c_0)$, for any scalar value $x$ and index $n$, we have
    \[ 
    \text{ExtConv}(c, \delta_0(x)) \land
    \text{ExtFP}_1(c, \delta_0(x), n) \land
    \means{c}(\delta_0(x))[n] = 0 \implies
    \text{IntFP}_1(c, \delta_0(x), n+1) \]
\end{lemma}

This lemma gives the concrete upper bound $n+1$ on the iteration at which IntFP holds.
It also allows us to prove $\mathsf{datalogQuery}$ is Conv-complete.
\begin{lemma}\label{thm:datalog-query-convcomplete}
    If $c_0$ is a Conv-complete regular circuit, then $\mathsf{datalogQuery}(c_0)$ is Conv-complete.
\end{lemma}

We have proved similar lemmas to establish the same iteration bound $n+1$ for $\mathsf{whileBody}$ and the Conv-completeness of $\mathsf{whileQuery}$.
These lemmas allow us to prove that regular circuits are Conv-complete by induction on the structure of regular circuits.
\begin{theorem}[Regular Circuits are Conv-complete]\label{thm:regular-convcomplete}
    Any regular circuit is Conv-complete.
\end{theorem}

From the proof above, we can see that the class of regular circuits is highly extensible. We can freely add more recursive circuit constructs similar to $\mathsf{datalogBody}$ and $\mathsf{whileBody}$ while preserving Conv-completeness, as long as we can establish the variant of FP-completeness in Thm.~\ref{thm:datalog-body-fpcomplete} for the new constructors.  

\subsection{Conv-complete Circuits Are Closed under Incremental Optimization}\label{sec:inc-opt}

As we can see in our example, the standard incremental optimization of DBSP can be represented by two syntactic transformations applied to a non-incremental circuit: $\mathsf{plift}$ to push the lifting operator onto individual nodes and $\mathsf{incOpt}$ to apply compositional incremental rules. In this section, we aim to show that $\mathsf{plift}$ and $\mathsf{incOpt}$ preserve Conv-completeness. To do this, we establish their correctness theorems to show how they transform the denotation, the IntFP, and the IntConv of the original circuit. The IntFP results track convergence bounds through these transformations; they are semantic guarantees rather than measurements of runtime cost.

To formally state the correctness of $\mathsf{plift}$, we introduce the notion of \emph{strong semantic equivalence} that requires the equivalence of both external denotation and internal convergence.
\begin{definition}[Strong Semantic Equivalence]
    For any circuit $c_1$ and $c_2$, we define the strong semantic equivalence $c_1 \simeq c_2$ as follows:
    \[\forall x,\, \means{c_1}(x) = \means{c_2}(x) \land (\text{IntConv}(c_1, x) \iff \text{IntConv}(c_2, x)) \]
    Note that this implicitly implies that $\text{ExtConv}(c_1, x) \iff \text{ExtConv}(c_2, x)$ because $\means{c}$ is partial.
\end{definition}

Based on this definition, we have re-established important equivalence lemmas in the DBSP theory, such as the chain rule $\Delta(c_1 \seq c_2) \simeq \Delta(c_1) \seq \Delta(c_2)$ and the lifting cycle rule $\mathsf{lifting}(\mathsf{loop}(c)) \simeq \mathsf{\lift{loop}}(\mathsf{lifting}(c))$. Using these, we prove the correctness of $\mathsf{plift}$ as follows. It is a straightforward induction on the structure of $c$.

\begin{theorem}[Correctness of $\mathsf{plift}$]\label{thm:pushLifting-correctness}
    For any level-1 circuit $c$,
    \begin{enumerate}
        \item (Denotation and IntConv) $\mathsf{lifting}(c) \simeq \mathsf{plift}(c)$.
        \item ($\text{IntFP}_1$) For any $x$ and $n$, $\text{ExtFP}_1(\mathsf{plift}(c), x, n) \implies \text{IntFP}_1(\mathsf{plift}(c), x, n)$.
        \item ($\text{IntFP}_2$) For any $x$ and $m,n \in \N$, $\text{IntFP}_1(c, x[m], n) \implies \text{IntFP}_2(\mathsf{plift}(c), x, m, n)$.
    \end{enumerate}
\end{theorem}

This allows us to prove that $\mathsf{plift}$ preserves both Conv-completeness and FP-completeness.
\begin{theorem}[$\mathsf{plift}$ preserves Conv-completeness and FP-completeness]\label{thm:pushLifting-convcomplete}
    For any level-1 circuit $c$, if $c$ is Conv-complete, then $\mathsf{plift}(c)$ is Conv-complete. If $c$ is FP-complete, then $\mathsf{plift}(c)$ is FP-complete.
\end{theorem}

For $\mathsf{incOpt}$, the situation is more tricky. In particular, unlike $\mathsf{lifting}(c)$ and $\mathsf{plift}(c)$, $\Delta(c)$ and $\mathsf{incOpt}(c)$ may have different IntConv behavior; just consider $\mathsf{node}(\mathsf{id})$ and $\Delta(\mathsf{node}(\mathsf{id}))$. This motivates us to introduce \emph{weak semantic equivalence}, which only requires the equivalence of external denotation but not internal convergence.
\begin{definition}[Weak Semantic Equivalence]
    For any circuit $c_1$ and $c_2$, we define the weak semantic equivalence $c_1 \approx c_2$ as $\forall x, \means{c_1}(x) = \means{c_2}(x)$. Note that this implies that $\text{ExtConv}(c_1, x) \iff \text{ExtConv}(c_2, x)$ since $\means{c}$ is a partial function.
\end{definition}

Besides denotation, how $\mathsf{incOpt}$ transforms IntFP is also tricky.
To see this, consider the simplest circuit $\mathsf{id}$, whose optimized incremental form is itself. Although the two forms are the same circuit, their corresponding input is different, because the input $x$ to the original form corresponds to the input $\D(x)$ to the incremental form. Because $x$ and $\D(x)$ have different FixAfter indices in general, the two forms have different IntFP in general.
This observation suggests that the formulation of the IntFP-preservation property that we desire for $\mathsf{incOpt}$ will be non-trivial.


An important observation is that although $x$ and $\D(x)$ may have different FixAfter indices, they do have a connection as follows:
\begin{proposition}
    For any stream $s$ and $n \in \N$, $\text{FixAfter}_1(s, n) \iff \text{ZeroAfter}_1(\D(s), n + 1)$,
    and for any nested stream $s$ and $\mathbf{b} \in \stream{\N}$, 
    $\text{FixAfter}_2(s, \mathbf{b}) \implies \text{ZeroAfter}_2(\D(s), \lift{\max}(\mathbf{b}, z^{-1}(\mathbf{b})))$.
\end{proposition}

Note that ZeroAfter implies FixAfter. This inspires us to raise a natural question: \emph{since input streams of $c$ and $\mathsf{incOpt}(c)$ have such FixAfter index relation, will all other intermediate streams share the same relation?} If so, as a result, IntConv should also be preserved. We formalize this idea as the following \emph{preservation} relation. 
\begin{definition}[IntFP and IntConv Preservation Relation]
    For any circuit $c_1$ and $c_2$, we say that $c_1$'s $\text{IntFP}_1$ is preserved by $c_2$, denoted by $c_1 \preserve_1 c_2$, if 
    \[\forall x, n,\, (\text{ExtConv}(c_1, x) \land \text{IntFP}_1(c_1, x, n)) \implies \text{IntFP}_1(c_2, \D(x), n + 1)\]
    For level-2 circuit $c_1$ and $c_2$, we say that $c_1$'s $\text{IntFP}_2$ is preserved by $c_2$, denoted by $c_1 \preserve_2 c_2$, if
    \[\forall x, \mathbf{b},\, (\text{ExtConv}(c_1, x) \land \text{IntFP}_2(c_1, x, \mathbf{b})) \implies \text{IntFP}_2(c_2, \D(x), \lift{\max}(\mathbf{b}, z^{-1}(\mathbf{b})))\]
    For any circuit $c_1$ and $c_2$, we say that $c_1$'s IntConv is preserved by $c_2$, denoted by $c_1 \preserve_{C} c_2$, if
    \[\forall x, \text{IntConv}(c_1, x) \implies \text{IntConv}(c_2, \D(x))\]
\end{definition}

These allow us to formally define what we expect from a ``good'' incremental optimization:
\begin{definition}[Sound Incremental Form]
    For a circuit $c_1$ and $c_2$, we say that $c_2$ is a \emph{sound incremental form} of $c_1$ if they satisfy the following four properties:
    \begin{enumerate}
        \item (Denotation) $\Delta(c_1) \approx c_2$.
        \item ($\text{IntFP}_1$) $c_1 \preserve_1 c_2$.
        \item ($\text{IntFP}_2$) $c_1 \preserve_2 c_2$ if $c_1$ is a level-2 circuit.
        \item (IntConv) $c_1 \preserve_{C} c_2$.
    \end{enumerate}
\end{definition}

Another important challenge is that $\mathsf{incOpt}$ is a flexible and extensible algorithm, as different primitive nodes can have different and custom incremental optimization. We deal with this via a uniform assumption that $\mathsf{incOpt}$ transforms any primitive node into a sound incremental form. This is a strong but still practical assumption, as we have established that the following common incremental forms are all sound:
\begin{proposition}[Common Sound Incremental Forms]\label{thm:common-sound-incremental-forms}
    \;
    \begin{enumerate}
        \item (Unoptimized) For any circuit $c$, $\Delta(c)$ is a sound incremental form of $c$.
        \item (Linear) For any linear function $f$ and $c = \mathsf{node}(f)$, $c$ is a sound incremental form of $c$.
        \item (Bilinear) For any bilinear function $f$ and $c = \mathsf{node}(f)$, the DBSP's bilinear incremental optimization~\cite{budiu2025dbsp} of $c$, denoted $\mathsf{bilinearOpt}(c)$, is a sound incremental form of $c$.
    \end{enumerate}
\end{proposition} 
The unoptimized incremental form provides a general sound fallback for any primitive node, ensuring the robustness of $\mathsf{incOpt}$. From now on, we may assume that $\mathsf{incOpt}$ transforms all its primitive nodes into a sound incremental form, and we can finally establish the correctness of $\mathsf{incOpt}$ as follows:
\begin{theorem}[Correctness of $\mathsf{incOpt}$]\label{thm:incOpt-correctness}
    For any circuit $c$, $\mathsf{incOpt}(c)$ is a sound incremental form of $c$.
\end{theorem}

This allows us to show that $\mathsf{incOpt}$ preserves Conv-completeness.
\begin{theorem}[$\mathsf{incOpt}$ preserves Conv-completeness]\label{thm:incOpt-convcomplete}
    For any circuit $c$, if $c$ is Conv-complete, then $\mathsf{incOpt}(c)$ is Conv-complete.
\end{theorem}

As a result, we can prove that the circuit implementing the delta-of-deltas strategy as shown in Fig.~\ref{fig:delta-of-deltas-dbsp} is Conv-complete.
\begin{proposition}
   If $c_0$ is a regular circuit, then $\mathsf{optDatalogQuery}(c_0)$ (Equation~\ref{eq:opt-datalog-query}) is Conv-complete.
\end{proposition}
